% Part: normal-modal-logic
% Chapter: frame-definability
% Section: properties-accessibility

\documentclass[../../../include/open-logic-section]{subfiles}

\begin{document}

\olfileid{nml}{frd}{acc}

\olsection{Sipat-Sipat Relasi Aksesibilitas}

Akeh !!{formula}s modal kang nyata-nyata dadi ciri kanggo sipat
relasi aksesibilitas kang prasaja, malah wis lumrah dikenal. Ing arah
siji, iki ateges yen saben model kang nduweni sawijining sipat bakal
ndadekake !!{formula} kang cocog (lan kabeh instansi substitusine)
bener. Kita miwiti saka limang tuladha klasik jinis relasi
aksesibilitas lan formula kang kabenerane dijamin dening relasi kasebut.

\begin{thm}\ollabel{thm:soundschemas}
  Upamane $\mModel{M} = \tuple{W, R, V}$ sawijining model. Yen $R$
  nduweni sipat ing sisih kiwa \olref{tab:five}, saben instansi saka
  !!{formula} ing sisih tengen bener ing~$\mModel{M}$.
\end{thm}

\begin{table}[t]
    \begin{tabular}{| l || l |}
      \hline
      {\emph{Yen $R$ iku \dots}} & {\emph{mula \dots bener ing~$\mModel{M}$:}} \\
      \hline \hline
      \emph{serial}: $\forall u \exists v Ruv$ & \hbox to.43\textwidth
           {$\Box p \lif \Diamond p$ \hfill (\Ax{D})} \\
      \hline
      \emph{refleksif}: $\forall w Rww$  
      & \hbox to.43\textwidth
          {$\Box p \lif p$ \hfill (\Ax{T})} \\
      \hline
      \emph{simetris}: &  \hbox to .43\textwidth
           {$p \lif \Box\Diamond p$ \hfill (\Ax{B})} \\
           $\forall u\forall v(Ruv \lif Rvu)$ & \\
      \hline
      \emph{transitif}: & \hbox to.43\textwidth
           {$\Box p \lif \Box \Box p$ \hfill (\Ax{4})} \\
      $\forall u \forall v \forall w ((Ruv \land Rvw) \lif Ruw)$ & \\
      \hline 
      \emph{Euklidean}: & \hbox to .43\textwidth
        {$\Diamond p \lif \Box\Diamond p$ \hfill (\Ax{5})} \\
      $\forall w \forall u \forall v ((Rwu \land Rwv) \lif Ruv)$ &\\
      \hline
    \end{tabular}
    \caption{Lima fakta korespondensi.}
    \ollabel{tab:five}
  \end{table}

\begin{proof}
  Iki kasus kanggo \Ax{B}: kanggo nuduhake yen skema iki bener ing
  sawijining model, kita kudu nuduhake yen kabeh instansine bener ing
  kabeh donya ing model kasebut. Mula pilih $!A \lif \Box\Diamond !A$
  minangka salah sawijining instansi \Ax{B}, lan pilih $w \in W$
  minangka donya sebarang. Upamane anteseden $!A$ bener ing $w$;
  kita arep nuduhake yen $\Box \Diamond !A$ bener ing $w$. Dadi kita
  kudu nuduhake yen $\Diamond !A$ bener ing kabeh $w'$ kang bisa diakses
  saka $w$. Saiki, kanggo saben $w'$ kang nyukupi $Rww'$, saka
  hipotesis simetri kita uga duwe $Rw'w$ (delengen
  \olref{fig:Bsymm}). Amarga $\mSat{M}{!A}[w]$, kita duwe
  $\mSat{M}{\Diamond !A}[w']$. Amarga $w'$ dipilih sebarang kanthi
  $Rww'$, kita duwe $\mSat{M}{\Box\Diamond!A}[w]$.

Kasus-kasus liyane kita pasrahake minangka gladhen.
\end{proof}

\begin{prob}
  Rampungna bukti \olref[nml][frd][acc]{thm:soundschemas}.
\end{prob}

\begin{figure}
  \begin{center}
    \begin{tikzpicture}[modal]
      \node[world] (w1) [label={below:$\mSat{{}}{\formula{A}}$\\
          $\mSat{{}}{\Box\Diamond \formula{A}}$}] {$w$}; 
      \node[world] (w2) [label=below:{$\mSat{{}}{\Diamond \formula{A}}$},
        right=of w1] {$w'$}; 
      \draw[->, bend left] (w1) to (w2); 
      \draw[->, bend left] (w2) to (w1); 
    \end{tikzpicture}
  \end{center}
\caption{Argumentasi saka simetri.}
\ollabel{fig:Bsymm}
\end{figure}

Gatekna yen implikasi kosok baline \olref{thm:soundschemas} ora
bener: yen sawijining model nyukupi skema, relasi aksesibilitas model
iku durung mesthi nduweni sipat kang cocog. Ing kasus \Ax{T} lan model
refleksif, gampang menehi tuladha model kang \Ax{T} dhewe gagal:
pilih $W = \{w\}$, $R = \emptyset$, lan
$V(p) = \emptyset$. Dadi $R$ ora refleksif, nanging $\mSat{M}{\Box
  p}[w]$ lan $\mSat/{M}{p}[w]$. Nanging ing kene mung siji instansi
\Ax{T} kang gagal ing~$\mModel{M}$; instansi liyane, upamane
$\Box \lnot p \lif \lnot p$, bener. Luwih angel menehi tuladha kang
\emph{saben instansi substitusi} saka~\Ax{T} bener
ing~$\mModel{M}$ nanging $\mModel{M}$ ora refleksif. Nanging model kaya
mangkono pancen ana:

\begin{prop}\ollabel{prop:reflexive}
  Upamane $\mModel{M} = \tuple{W, R, V}$ sawijining model kanthi $W = \{u, v
  \}$, ing ngendi donya $u$ lan $v$ direlasikake dening $R = \{(u,v),(v,u)\}$:
  yaiku $Ruv$ lan $Rvu$, tanpa gelang dhiri. Upamane kanggo saben $p$:
  $u \in V(p) \Leftrightarrow v
  \in V(p)$. Mula:
  \begin{enumerate}
  \item Kanggo saben $!A$: $\mSat{M}{!A}[u]$ yen lan mung yen
    $\mSat{M}{!A}[v]$ (gunakna induksi marang $!A$).
  \item Saben instansi saka~\Ax{T} bener ing $\mModel{M}$.
  \end{enumerate}
  Amarga $\mModel{M}$ ora refleksif (malah
  \emph{irrefleksif}), implikasi kosok baline \olref{thm:soundschemas}
  gagal kanggo \Ax{T} (argumentasi kang padha bisa diwenehake kanggo
  saweneh---nanging ora kabeh---skema liyane ing
  \olref{thm:soundschemas}).
\end{prop}

\begin{prob}
  Buktekna pratelan ing \olref[nml][frd][acc]{prop:reflexive}.
\end{prob}

Senadyan kita bakal nggatekake limang !!{formula}s klasik \Ax{D},
\Ax{T}, \Ax{B}, \Ax{4}, lan \Ax{5}, kita cathet ing
\olref{tab:anotherfive} sawetara sipat relasi aksesibilitas liyane.
Relasi aksesibilitas~$R$ iku fungsional parsial yen saka saben donya
paling akeh siji donya kang bisa diakses. Yen saka saben donya pas siji
donya kang bisa diakses, relasi iku diarani fungsional. (Dadi relasi
fungsional persis relasi kang serial lan fungsional parsial.) Relasi
iki diarani ``fungsional'' amarga relasi aksesibilitas mau tumindak
kaya fungsi (parsial). Sawijining relasi diarani padhet lemah yen
saben $Ruv$ ana donya~$w$ ``ing antarane'' $u$ lan~$v$. Dadi relasi
padhet lemah, ing sawijining pangertèn, kosok baline relasi transitif:
ing relasi transitif, yen $v$ bisa digayuh saka~$u$ liwat~$w$, $v$
uga bisa digayuh langsung saka $u$; ing relasi padhet lemah, yen $v$
bisa digayuh langsung saka~$u$, donya iku uga bisa digayuh kanthi liwat
sawijining~$w$. Sawijining relasi diarani kaarahake lemah yen, saben
donya $u$ lan~$v$ bisa digayuh saka sawijining donya~$w$, ana siji
donya~$t$ kang bisa digayuh saka $u$ lan uga saka~$v$---iki kadhang
diarani ``sipat inten'' utawa ``konfluensi.''

\begin{table}[t]
    \begin{tabular}{| p{.50\textwidth} || p{.43\textwidth} |}
      \hline
      {\emph{Yen $R$ iku \dots}} & {\emph{mula \dots bener ing~$\mModel{M}$:}} \\
      \hline \hline
      \emph{fungsional parsial}: & 
      \multirow{2}{*}{$\Diamond p \lif \Box p$} \\
      $ \forall w \forall u \forall v ((Rwu \land Rwv) \lif u =v )$ & \\
      \hline
      \multirow{2}{*}{\emph{fungsional}: $\forall w \exists v \forall u(Rwu \liff \eq[u][v]) $} 
      & \multirow{2}{*}{$\Diamond p \liff \Box p$} \\
      & \\
      \hline
      \emph{padhet lemah}: &  \multirow{2}{*}{$\Box \Box p \lif
        \Box p$} \\
      $\forall u\forall v(Ruv \lif \exists w(Ruw \land Rwv))$ & \\
      \hline
      \emph{kasambung lemah}: &
      \multirow{3}{*}{\hbox to.43\textwidth{$
        \begin{array}{@{}l@{}}
          \Box ((p \land \Box p) \lif q) \lor {} \\
          \qquad \Box ((q \land \Box q) \lif p)
        \end{array}$ \hfill (\Ax{L})}
      } \\
      $\forall w \forall u \forall v ((Rwu \land Rwv) \lif {}$ &\\
      \qquad $(Ruv \lor u=v \lor Rvu))$ & \\
      \hline 
      \emph{kaarahake lemah}: & \multirow{3}{*}{\hbox to .43\textwidth
        {$\Diamond\Box p \lif \Box\Diamond p$\hfill (\Ax{G})}} \\
      $\forall w \forall u \forall v ((Rwu \land Rwv) \lif {}$ & \\
      \qquad $\exists t (Rut \land Rvt))$ & \\
      \hline
    \end{tabular}
    \caption{Lima fakta korespondensi liyane.}
    \ollabel{tab:anotherfive}
  \end{table}

\begin{prob}
  Upamane $\mModel{M} = \tuple{W, R, V}$ sawijining model. Tuduhna yen
  $R$ nyukupi sipat ing sisih kiwa
  \olref[nml][frd][acc]{tab:anotherfive}, saben instansi formula kang
  cocog ing sisih tengen bener ing~$\mModel{M}$.
\end{prob}

\end{document}
